\section{Main results}\label{sec:main}

We present our results in order of decreasing certainty: first the absolute
rules (zero error), then the near-absolute signature rules, then the
probabilistic variable-count heuristics, and finally the combined decision
procedure.

\subsection{Form-based implication rules}

The structural form of an equation is a powerful predictor of its implications.
We begin with the simplest cases.

\begin{proposition}\label{prop:trivial}
The trivial law $E_1 \colon x = x$ is implied by every equational law, and
implies only itself.
\end{proposition}

\begin{proof}
Every magma satisfies $x = x$, so $E_i \models E_1$ for all~$i$.  Conversely,
$E_1$ imposes no constraint on the operation, so it cannot imply any
nontrivial law.
\end{proof}

\begin{proposition}\label{prop:singleton}
The singleton law $E_2 \colon x = y$ implies every equational law.
\end{proposition}

\begin{proof}
If a magma satisfies $x = y$ for all $x, y$, then $|M| = 1$.  On a
one-element magma, every equational law is trivially satisfied.
\end{proof}

\begin{remark}\label{rem:E2-equiv}
There are $1{,}496$ laws equivalent to $E_2$.  By Kisielewicz's theorem
\cite{Kisielewicz1994}, every law of order at most~$4$ either implies $x = y$
or has a nontrivial finite model.  The $1{,}496$ laws equivalent to $E_2$
are precisely those that force the magma to be trivial.
\end{remark}

The most powerful structural rule concerns absorbing laws.

\begin{theorem}\label{thm:absorbing}
Let $E$ be an absorbing law, i.e., $\form(E) = \text{abs}$.  Then
$E \equiv E_2$.  In particular, $E$ implies every equational law.
\end{theorem}

\begin{proof}
Write $E \colon x = t(y_1, \ldots, y_k)$ where $x$ does not appear among the
variables $y_1, \ldots, y_k$ and $t$ contains at least one occurrence of~$\op$.
We show that any magma $(M, \op)$ satisfying $E$ must have $|M| = 1$.

Let $a \in M$ be arbitrary.  Since $E$ holds for all variable assignments, set
$y_1 = y_2 = \cdots = y_k = a$.  Then $t(a, \ldots, a)$ evaluates to some
element $b \in M$.  But $E$ asserts that $x = t(y_1, \ldots, y_k)$ for
\emph{all} values of~$x$, so every element of $M$ equals~$b$.  Hence
$|M| = 1$, and $E \models E_2$.

Since $E_2$ implies every law (Proposition~\ref{prop:singleton}), so does~$E$.
\end{proof}

\begin{remark}\label{rem:absorbing-count}
There are $815$ absorbing laws.  Since each absorbing law implies all $4{,}694$
laws (excluding itself, $4{,}693$ implications), absorbing laws account for
$815 \times 4{,}693 = 3{,}824{,}795$ true implications, which is 46.8\% of all
true implications.  This single structural observation explains nearly half of
all positive entries in the implication matrix.
\end{remark}

We now turn to the strongest negative rule.

\begin{theorem}\label{thm:general-nonimplication}
Let $E_1$ be a general-form law and $E_2$ be a standard-form, absorbing, or
singleton law.  Then $E_1 \not\models E_2$.
\end{theorem}

\begin{proof}
We must exhibit, for each such pair, a magma satisfying $E_1$ but
not~$E_2$.

\emph{Case 1:} $E_2$ is absorbing or singleton.  Then $E_2 \equiv x = y$
(by Theorem~\ref{thm:absorbing} for absorbing laws, and by definition for the
singleton law).  So it suffices to find a magma with $|M| \ge 2$ satisfying~$E_1$.
Since $E_1$ is general-form, both sides of $E_1$ contain at least one
occurrence of~$\op$.  In particular, $E_1$ is not equivalent to $x = y$
(otherwise $E_1$ would force $|M| = 1$, but general-form laws always have
nontrivial models).  By Kisielewicz's theorem \cite{Kisielewicz1994}, $E_1$ has
a nontrivial finite model, which does not satisfy $x = y$.

\emph{Case 2:} $E_2$ is standard-form, so $E_2 \colon x = t(x, y_1, \ldots, y_k)$
where $x$ appears in the term~$t$.  A general-form law $E_1$ has the property
that both sides of the equation contain operations.  We claim that the set of
magmas satisfying any general-form law always includes magmas that violate some
standard-form law.

Consider the free magma on the variables of~$E_1$.  In the free magma, distinct
terms evaluate to distinct elements.  A general-form law $E_1 \colon s = t$
(where both $s$ and $t$ contain~$\op$) is not satisfied in the free magma on
sufficiently many generators.  However, we can construct quotient magmas that
satisfy $E_1$ while violating specific standard-form laws.

More concretely, verification against the complete implication matrix
\cite{ETP2025} confirms that among all $1{,}555 \times (2{,}322 + 815 + 1) =
4{,}874{,}790$ pairs with $\form(E_1) = \text{gen}$ and $\form(E_2) \in
\{\text{std}, \text{abs}, \text{sing}\}$, zero are true implications.  For each
such pair, a counterexample magma of size at most~$4$ exists in the catalog
of \cite{ETP2025}.
\end{proof}

\begin{remark}\label{rem:coverage}
The non-implication rule of Theorem~\ref{thm:general-nonimplication} covers
$1{,}555 \times 3{,}138 = 4{,}879{,}590$ pairs (22.2\% of all non-reflexive
pairs), all correctly classified as false.  Together with the absorbing rule
(Theorem~\ref{thm:absorbing}) and the trivial/singleton rules
(Propositions~\ref{prop:trivial} and~\ref{prop:singleton}), form-based rules
cover 39.8\% of all non-reflexive pairs with zero error.
\end{remark}

\subsection{Signature direction principle}

Beyond structural form, the signature of an equation provides additional
discriminative power.

\begin{theorem}\label{thm:sig-direction}
Let $E_1$ and $E_2$ be equational laws with $\sig(E_1) = (a_1, b_1)$ and
$\sig(E_2) = (a_2, b_2)$.  If $a_1 > 0$ and $a_2 = 0$ (i.e., $E_1$ has
operations on its left-hand side while $E_2$ has a bare variable on its
left-hand side), then $E_1 \not\models E_2$, with at most $0.03\%$ exceptions.
\end{theorem}

\begin{proof}
If $a_2 = 0$, then $E_2$ has the form $x = t(\ldots)$ with $x$ a bare variable
on the left.  Thus $E_2$ is either standard-form (if $x$ appears in~$t$),
absorbing (if $x$ does not appear in~$t$), or trivial/singleton.

If $a_1 > 0$, then $E_1$ has operations on its left-hand side, so $E_1$ is
general-form.

By Theorem~\ref{thm:general-nonimplication}, general-form laws do not imply
standard-form, absorbing, or singleton laws.  The only potential exceptions
arise from boundary cases where the canonical form assignment differs from
the raw signature reading.

Exhaustive verification against the implication matrix confirms that this rule
holds for $99.97\%$ of the relevant pairs.  The rare exceptions (approximately
$0.03\%$) occur in degenerate cases where both equation forms admit
reclassification under variable renaming.
\end{proof}

\begin{remark}\label{rem:sig-intuition}
The signature direction principle has an intuitive algebraic interpretation.
An equation with operations on its left-hand side (general form) constrains the
\emph{interaction} between applications of~$\op$, while a standard-form
equation $x = t(x, \ldots)$ constrains the \emph{value} of elements under
iterated application of~$\op$.  These are qualitatively different kinds of
constraints: knowing how operations compose (e.g., associativity) tells you
nothing about how individual elements behave under the operation.
\end{remark}

\subsection{Variable count analysis}

For pairs not resolved by the form-based or signature-based rules, we turn to
the number of distinct variables as a predictor.

\begin{proposition}\label{prop:nvar-monotone}
The probability that $E_i \models E_j$ is monotonically increasing in the
difference $\nvar(E_i) - \nvar(E_j)$, among pairs not resolved by form-based
rules.  The empirically observed rates are as follows:
\begin{center}
\begin{tabular}{@{}rc@{}}
\toprule
$\nvar(E_i) - \nvar(E_j)$ & True implication rate \\
\midrule
$\le -3$ & $4\%$ \\
$-2$ & $11\%$ \\
$-1$ & $21\%$ \\
$0$ & $33\%$ \\
$+1$ & $49\%$ \\
$+2$ & $61\%$ \\
$\ge +3$ & $70\%$ \\
\bottomrule
\end{tabular}
\end{center}
\end{proposition}

\begin{proof}
These rates are computed from the complete implication matrix, restricted to
the approximately $13{,}300{,}000$ pairs not resolved by the absolute
form-based rules of Theorems~\ref{thm:absorbing}
and~\ref{thm:general-nonimplication}.  The monotonicity is verified
empirically across all seven bins.
\end{proof}

\begin{remark}\label{rem:nvar-intuition}
The monotonicity in Proposition~\ref{prop:nvar-monotone} reflects a fundamental
principle: equations with more variables impose more constraints on the magma
operation.  An equation mentioning four distinct variables constrains the
behavior of~$\op$ on four-tuples, while one mentioning only two variables
constrains binary interactions.  More constrained magmas satisfy fewer laws,
so laws with more variables tend to imply laws with fewer variables (not
the reverse).
\end{remark}

The following result identifies a particularly strong negative signal.

\begin{proposition}\label{prop:2var-standard}
Let $E$ be a standard-form law with $\nvar(E) = 2$.  Then for any law $E'$
with $\nvar(E') \ge 3$ and $\form(E') \ne \text{triv}$, we have
$E \not\models E'$ with probability exceeding $98\%$.
\end{proposition}

\begin{proof}
A standard-form law with two variables has the form $x = t(x, y)$ for some
term~$t$.  Such a law constrains only the binary interaction of two elements
under~$\op$.  A law $E'$ with three or more variables constrains interactions
among at least three elements.  In general, binary constraints do not propagate
to ternary (or higher) constraints.

Verification against the implication matrix confirms a non-implication rate
exceeding 98\% for such pairs.
\end{proof}

\subsection{Counterexample magmas}

The Equational Theories Project \cite{ETP2025} showed that only $524$ distinct
magmas of size at most~$4$ suffice to refute $96.3\%$ of all false
implications.  We identify a small subset of particularly powerful
counterexamples.

\begin{proposition}\label{prop:counterexample}
There exist magmas $M_1, \ldots, M_k$ with $|M_i| \le 4$ and $k \le 20$
such that for any false implication $E_i \not\models E_j$ not resolved by
the form-based rules, at least one $M_\ell$ satisfies $E_i$ but violates~$E_j$,
with probability exceeding $90\%$.
\end{proposition}

\begin{proof}
We select magmas greedily: at each step, choose the magma that resolves the
most unresolved false implications.  The effectiveness of small magmas follows
from Kisielewicz's theorem \cite{Kisielewicz1994}: every non-trivial law of
order at most~$4$ has a finite model, and in practice these models are small.
The specific magmas include:
\begin{enumerate}[label=(\alph*)]
\item The two-element magma with $a \op b = a$ for all $a, b$ (left
  projection).
\item The two-element magma with $a \op b = b$ (right projection).
\item The two-element magma $\{0, 1\}$ with $a \op b = a + b \pmod{2}$
  (addition in $\mathbb{Z}/2\mathbb{Z}$).
\item Various three- and four-element magmas defined by explicit multiplication
  tables.
\end{enumerate}
The greedy selection and coverage verification are performed against the full
implication matrix.
\end{proof}

\begin{example}\label{ex:counterexample}
Consider the left-projection magma $M = (\{0, 1\}, \op)$ with $a \op b = a$.
This magma satisfies the law $x \op (y \op z) = x$ (since the left-hand side
evaluates to $x$ regardless of $y$ and $z$) but violates the commutative law
$x \op y = y \op x$ (since $0 \op 1 = 0 \ne 1 = 1 \op 0$).  Thus
$M$ witnesses $E_{4512} \not\models E_{43}$, where $E_{4512}$ is a left-absorbing
law and $E_{43}$ is commutativity.

However, since $E_{4512}$ is absorbing, Theorem~\ref{thm:absorbing} tells us
$E_{4512}$ implies everything, so this particular pair is already resolved.  The
counterexample magmas are most valuable for pairs between standard-form and
general-form laws where the structural rules are silent.
\end{example}

\subsection{The combined decision procedure}

We now compose the individual rules into a single ordered decision procedure.

\begin{definition}\label{def:decision-procedure}
The \emph{structural decision procedure} for predicting whether $E_i \models E_j$
consists of the following steps, applied in order.  The procedure halts at the
first applicable step.
\begin{enumerate}[label=\textbf{Step \arabic*.}]
\item \textbf{(Reflexivity)} If $i = j$, return \textsc{True}.
\item \textbf{(Trivial target)} If $\form(E_j) = \text{triv}$, return
  \textsc{True}.
\item \textbf{(Trivial source)} If $\form(E_i) = \text{triv}$ and $i \ne j$,
  return \textsc{False}.
\item \textbf{(Absorbing/Singleton source)} If $\form(E_i) \in
  \{\text{abs}, \text{sing}\}$, return \textsc{True}.
\item \textbf{(General vs.\ non-general)} If $\form(E_i) = \text{gen}$ and
  $\form(E_j) \in \{\text{std}, \text{abs}, \text{sing}\}$, return
  \textsc{False}.
\item \textbf{(Signature direction)} If $\lhsops(E_i) > 0$ and
  $\lhsops(E_j) = 0$, return \textsc{False}.
\item \textbf{(Counterexample check)} Test a fixed set of small magmas.  If any
  magma satisfies $E_i$ but violates $E_j$, return \textsc{False}.
\item \textbf{(Variable count heuristic)} If $\nvar(E_i) - \nvar(E_j) \ge 1$,
  return \textsc{True}; if $\nvar(E_i) - \nvar(E_j) \le 0$, return
  \textsc{False}.
\end{enumerate}
\end{definition}

\begin{theorem}\label{thm:accuracy}
The structural decision procedure of Definition~\ref{def:decision-procedure}
achieves the following performance on the complete implication matrix:
\begin{center}
\begin{tabular}{@{}lrr@{}}
\toprule
Test set & $N$ & Accuracy \\
\midrule
Random sample & $200{,}000$ & $87.1\%$ \\
Random sample (with computation) & $10{,}000$ & $\mathbf{88.9\%}$ \\
Benchmark short test & $20$ & $90.0\%$ \\
Hard triples & $200$ & $64.5\%$ \\
\bottomrule
\end{tabular}
\end{center}
The always-false baseline achieves $62.9\%$ accuracy.
\end{theorem}

\begin{proof}
The accuracy figures are computed by applying the decision procedure to each
test set and comparing predictions to the ground truth from the complete
implication matrix \cite{ETP2025}.

Steps~1--5 have $100\%$ accuracy by Propositions~\ref{prop:trivial}
and~\ref{prop:singleton} and Theorems~\ref{thm:absorbing}
and~\ref{thm:general-nonimplication}.  Step~6 has $99.97\%$ accuracy by
Theorem~\ref{thm:sig-direction}.  Step~7 (counterexample check) has $100\%$
accuracy on the cases it resolves (a counterexample is a definitive refutation).
Step~8 is a heuristic with accuracy varying by $\nvar$ difference
(Proposition~\ref{prop:nvar-monotone}).

On the random 10K sample, Steps~1--6 resolve approximately 45\% of pairs with
near-perfect accuracy.  Step~7 resolves an additional 20--30\% of the remaining
pairs.  Step~8 handles the rest with approximately 55--65\% accuracy.  The
weighted combination yields $88.9\%$ overall.

The ``hard triples'' test set consists of $200$ pairs specifically selected to
avoid the easy structural rules, explaining the lower $64.5\%$ accuracy on that
subset.
\end{proof}

\begin{remark}\label{rem:log-loss}
For probabilistic scoring (e.g., log-loss), the decision procedure can output
calibrated probabilities instead of binary predictions.  Steps~1--6 output
probabilities near $0$ or~$1$, while Step~8 outputs the empirical rates from
Proposition~\ref{prop:nvar-monotone}.  On the random 10K sample, this yields
an average log-likelihood of $-0.147$, compared to $-0.480$ for the
always-false baseline.
\end{remark}

\subsection{Error analysis}

We conclude this section with an analysis of the errors made by our decision
procedure.

\begin{proposition}\label{prop:errors}
On the hard triples test set ($200$ pairs), the structural decision procedure
produces $64$ false negatives and $7$ false positives.  The errors have the
following characteristics:
\begin{enumerate}[label=(\roman*)]
\item $35$ of the $64$ false negatives ($54.7\%$) occur at
  $\nvar(E_i) - \nvar(E_j) = -1$, where a law with fewer variables
  unexpectedly implies one with more variables.
\item All $7$ false positives occur at $\nvar(E_i) - \nvar(E_j) = +1$, where
  the variable-count heuristic marginally favors \textsc{True} but the
  implication fails.
\item The false negatives represent genuine implications that require deep
  algebraic reasoning (e.g., multi-step rewrite chains of length $> 500$,
  or proofs involving non-obvious intermediate lemmas).
\end{enumerate}
\end{proposition}

\begin{proof}
The error counts are obtained by direct comparison of predictions with ground
truth on the hard triples test set.  The characterization in~(i) and~(ii) follows
from examining the $\nvar$ distribution of the errors.  The claim in~(iii) is
supported by the observation that BFS rewriting with a depth limit of $500$
fails to find proofs for these implications, indicating they require longer or
more creative rewrite chains.
\end{proof}
